\section{Introduction}

\subsection{Queue Workers as Task Execution Systems}
% Object of analysis: task attempts inside queue-worker execution systems.
% Worker occupancy introduced briefly.
% Measurement-level mismatch introduced briefly.

Queue-based execution is a common design pattern for decoupling work production
from work execution. A producer submits a task instance to a queue, and one or
more workers later dequeue, reserve, execute, retry, acknowledge, or discard
that work. In small systems, this mechanism is often treated as a background-job
abstraction. In production systems, however, queue workers frequently become a
general execution layer for heterogeneous and operationally significant work:
data replication, cache rebuilds, database repair, media transformation,
external service calls, batch imports, validation, notification delivery,
maintenance jobs, and other forms of asynchronous processing.

This distinction matters because a queue worker is not merely a consumer of
messages. It is an execution boundary. Once a task attempt enters a worker, it
may consume CPU time, memory, disk I/O, network bandwidth, runtime allocations,
dependency capacity, database connections, leases, locks, semaphores, retry
budgets, and tenant-level concurrency slots. Some of these costs appear as
conventional hardware or runtime resource usage. Others are better understood
as worker-occupancy costs: the attempt holds a worker slot, a lease, a
connection, or another bounded execution resource for some duration. A task
attempt that consumes little CPU but occupies a worker slot for a long time may
reduce effective throughput and increase queueing delay more severely than a
shorter compute-dominant attempt.

The operational unit that profiling must explain is therefore not only the
worker process, container, or machine. It is the task attempt executed inside a
queue-worker boundary, together with the task kind, task-instance features, and
execution context that shaped that attempt. Section~\ref{sec:terminology-notation}
defines these terms more precisely. Informally, a task kind describes what work
is requested, a task instance is the concrete queued unit of that kind, and a
task attempt is one concrete execution of that instance.

However, the measurements easiest to collect are often not task-attempt-level
measurements. Operating systems and container runtimes expose CPU, memory,
I/O, and network counters for processes, cgroups, containers, or nodes. Worker
frameworks expose aggregate queue depth, worker utilization, task duration,
failure counts, and retry counts. Observability systems expose metrics, logs,
traces, and runtime profiles. These observations may show that a worker is
busy, that a queue is growing, or that a dependency is slow, but they do not by
themselves explain which task kind, task instance, or task attempt created the
pressure, which features were relevant, or which context factors amplified the
observed behavior.

This paper therefore treats queue workers as task execution systems rather
than passive queue consumers. Under this framing, task profiling becomes a
structured estimation problem: given a task kind, task-instance features, and
execution-context observations, what resource demand, worker occupancy,
uncertainty, and operational risk should the system expect for a task attempt,
and what profiling policy is justified? Different task kinds may require
different observations, which makes a single universal metric set insufficient.
The difficulty is deeper, however: even attempts of the same task kind may
produce substantially different costs under different execution contexts.

\subsection{Task Cost Is Conditional, Not Intrinsic}
% Same task kind and similar task-instance features may have different cost under different context.
% Intrinsic baseline + contextual amplification + residual uncertainty.

A natural first approximation is to assign a stable cost to a task kind, a
queue, or a payload-size bucket. Under this approximation, a
\texttt{replicate-partition} task with a one-gigabyte payload, an image
transformation task with a fixed input size, or a validation task for a small
document should have a predictable resource footprint. Resource-aware cluster
managers already reason about workloads through resource requests,
constraints, scheduling decisions, and observed usage~\cite{verma2015borg}.
However, prior work also shows that static or user-provided resource
assumptions may be insufficient for dynamic workloads. Quasar, for example,
motivates inferring resource requirements from application behavior and
performance constraints rather than relying only on user-specified
reservations~\cite{delimitrou2014quasar}. For queue-worker profiling, this
implies that task kind and declared size can provide a useful baseline, but
they should not be treated as the complete cost model.

Consider two task instances of the same replication task kind with the same
payload size, the same declared operation, and the same nominal priority. The
first instance is attempted against a healthy replica set with stable network
connectivity, low peer latency, and small replication lag. It completes in
roughly one minute. The second instance is attempted against a degraded replica
set with unstable network connectivity, peer delay, retransmissions, and
repeated retry or backoff intervals. It occupies a worker slot for fifty
minutes. The two instances are similar by task kind and declared task
features, but their attempts are not similar by operational behavior.

The difference is not explained by the intrinsic task description alone. The
declared intrinsic work may appear similar: read a range of data, transfer or
apply it to a peer, verify progress, and commit a replication checkpoint. The
observed behavior, however, is shaped by execution context. Network
instability can increase transfer time and retries. Peer degradation can
increase dependency wait time. Replication lag can correlate with a larger
reconciliation backlog. Disk pressure can delay reads or writes. Worker
pressure can increase local scheduling delay. The same baseline task kind can
therefore produce different resource demand and different worker occupancy
under different contexts. This view is consistent with production workload
studies that characterize task usage shapes and observe heterogeneous, dynamic
behavior in large compute clusters~\cite{zhang2011taskusage,reiss2012heterogeneity,cheng2018alibaba}.

This conditional behavior is not specific to replication. A transformation
task kind may exhibit compute-dominant behavior when record count and algorithm
mode dominate CPU time, but memory-pressure-sensitive behavior when allocation
volume or working-set size dominates. A task kind that calls external services
may be shaped by dependency latency, timeout policy, fan-out width,
connection-pool availability, and retry behavior. A database repair or cache
rebuild task kind may depend on hidden state such as segment fragmentation,
compaction debt, replica divergence, dirty data volume, or disk queue depth.
Prior workload-classification systems such as Paragon and Quasar also treat
workload behavior as something that must be classified with respect to
heterogeneity, interference, or performance constraints rather than inferred
from a name alone~\cite{delimitrou2013paragon,delimitrou2014quasar}.

Task profiling should therefore not assign a single intrinsic cost to a task
kind. A useful model must treat task cost as conditional: intrinsic task
features define a baseline, execution context can amplify or suppress that
baseline, and residual uncertainty remains when observed behavior is not
explained by known features. This residual uncertainty matters because rare
slow executions can dominate system-level behavior in large-scale distributed
systems~\cite{dean2013tail}. In other words, the profiling problem is not to
attach a universal cost label to a task kind, but to estimate
task-kind-specific behavior under a given execution context. Before selecting
metrics or instrumentation, the system must identify which task-instance
features and context factors are relevant for the task kind being profiled.

\subsection{Observations Are Not a Profiling Model}
% Metrics/traces/counters are useful observations.
% They do not define relevance, influence roles, double counting, feature mapping,
% or profiling strategy.

Metrics, logs, traces, counters, and runtime profiles are necessary inputs to
task profiling, but they are not a profiling model by themselves. Distributed
tracing systems such as Dapper demonstrate the value of tracing for
understanding large-scale distributed behavior, and modern telemetry systems
provide structured metrics, traces, and profiles for production
systems~\cite{sigelman2010dapper,opentelemetryMetricsDataModel}. These
mechanisms expose observations about execution. They do not, however, define
which observations are relevant for a particular task kind, which observations
should be interpreted as context factors, root-cause candidates, symptoms, or
outcomes, or how observations should be converted into estimates of resource
demand and worker occupancy.

The relevance problem appears immediately when different task kinds are
compared. A replication task kind may require observations about network
condition, peer health, replication lag, retries, and worker-slot time. A
transformation task kind may require input size, record count, CPU time,
allocation behavior, and memory pressure. A task kind that calls external
services may require dependency latency, timeout rate, retry count, and
connection-pool wait time. A simple validation task kind may require only
minimal accounting. A flat observation set shared by all task kinds either
misses important task-specific factors or collects unnecessary signals on the
hot path.

Raw observations also do not define influence structure. For example, network
degradation may increase transfer time, dependency latency, retries,
wall-clock duration, and worker-slot occupancy. Treating network throughput,
latency, retry count, and wall-clock duration as independent cost factors can
count the same underlying degradation multiple times. Conversely, treating only
the final wall-clock duration as the cost hides the mechanism that produced it.
A profiling model therefore needs to distinguish context factors,
root-cause candidates, symptoms, outcomes, derived signals, and residuals.

A further gap is feature mapping. A raw observation such as packet loss,
replica lag, connection-pool wait time, or allocation rate is not yet a model
input with stable meaning. It must be interpreted, normalized, and mapped into
a canonical feature such as network instability, replica health, dependency
capacity pressure, or runtime allocation pressure. Without this layer, two
domains may use different raw metrics for the same underlying influence, and
two similar metric names may represent different operational meanings. A task
profiling model must therefore specify how raw domain-specific observations
become canonical features used for estimation.

Finally, observations have a cost. Tracing, metric collection,
high-cardinality labels, runtime profiling, and telemetry export can add CPU,
memory, allocation, storage, and network overhead. Dapper was explicitly
designed around low overhead and application-level transparency, and
OpenTelemetry documentation similarly emphasizes that collecting more telemetry
can increase overhead and that high-cardinality metric tags can increase memory
usage~\cite{sigelman2010dapper,opentelemetryPerformance}. Sampling mechanisms,
including tail sampling, are therefore profiling-policy choices rather than
properties of observations themselves~\cite{opentelemetrySampling}. A profiling
system must decide what to observe, at what level of detail, for which task
kind or attempt, and under which overhead budget.

For these reasons, the problem addressed in this paper is not simply how to
collect more telemetry. The problem is how to interpret observations as part of
a task-kind-specific profiling model. The next subsection introduces influence
domains as the vocabulary for assigning semantic roles to observations and
Task Behavior Graphs as the structure for connecting those roles within a
task behavior model.

\subsection{From Influence Domains to Task Behavior Graphs}
% Influence domains provide vocabulary.
% Task Behavior Graphs connect domains for task kinds or task families.
% Task attempts are evaluated through the graph.
% Output: resource demand, occupancy, uncertainty, risk, profiling policy.

The preceding discussion suggests that task profiling requires more than
collecting additional observations. A profiling system must know what an
observation means for a given task kind, whether it describes the task
identity, intrinsic demand, execution context, resource demand, worker
occupancy, outcome, uncertainty, or risk, and whether it should affect resource
estimation, worker-occupancy estimation, uncertainty, or profiling policy. We
use \emph{influence domains} as the semantic vocabulary for making these
distinctions.

An influence domain groups model elements by their role in task behavior.
Intrinsic-demand features describe the work contained in the task instance,
such as payload size, record count, operation mode, or partition size.
Execution-context features describe the environment in which a task attempt
executes, such as network instability, dependency health, replica state, cache
warmth, disk pressure, or worker pressure. Resource-demand dimensions describe
active resource use, such as CPU time, memory peak, disk I/O, and network I/O.
Worker-occupancy dimensions describe bounded execution capacity held by the
attempt, such as worker-slot time, lease time, connection hold time, lock hold
time, or tenant-level concurrency time. Outcome, uncertainty, and risk domains
describe how the attempt completed, how predictable the estimate is, and what
consequences follow from estimation or control error.

Influence domains provide the vocabulary, but they do not by themselves define
a model. The model must also describe how domain-specific features influence
one another for a particular task kind or task family. This paper introduces
\emph{Task Behavior Graphs} for that purpose. A Task Behavior Graph represents
a task kind or task family as a typed graph whose nodes belong to influence
domains and whose edges encode modeled relationships such as baseline demand,
contextual amplification, additive overhead, derived symptoms, outcome links,
uncertainty indicators, and risk propagation. The graph is not assumed to prove
causality by itself; rather, it makes the model's influence assumptions
explicit so that they can be estimated, validated, calibrated, or revised.

For example, in a replication task kind, payload size may contribute to network
and disk demand, network instability may amplify retry behavior and
worker-slot time, and prolonged worker-slot time may increase queue-starvation
risk. In a transformation task kind, record count and algorithm mode may
instead dominate CPU time, allocation behavior, and memory pressure, while
network instability may be irrelevant or only weakly relevant. The graph
therefore does not encode a universal metric list. It encodes which features
and relationships are relevant for the task kind or family being modeled.

The graph structure is associated with a task kind or task family rather than
constructed independently for every task attempt. The task kind or family
provides the structure: relevant features, resource and occupancy dimensions,
influence roles, estimated sensitivities, and confidence in those relationships.
A task instance supplies intrinsic feature values, such as payload size or
record count. A task attempt supplies execution-context observations, attempt
state, worker state, and runtime observations. Under the model, graph
evaluation produces task-kind-specific estimates of resource demand, worker
occupancy, uncertainty, residuals, and operational risk. A separate
profiling-policy selector can then use these estimates to choose an
observation strategy, such as minimal accounting, detailed accounting, sampled
tracing, tail-aware profiling, diagnostic profiling, or isolated diagnostic
execution.

This framing also avoids the assumption that a profiling system can maintain a
globally complete metric list. Different domains may expose different raw
metrics for the same underlying influence, and different task kinds may require
different subsets of those influences. Instead of requiring every metric to be
known by the core model, Task Behavior Graphs operate over canonical features.
Domain-specific metrics can be mapped into these features through task behavior
manifests and domain packs, while unexplained behavior remains visible as
residual error and reduced confidence. The result is a task-kind-specific
profiling model: raw observations are interpreted through influence domains,
connected through a behavior graph, and used to select the amount and kind of
profiling justified for the task attempt.

\subsection{Contributions and Paper Organization}
% Contributions and roadmap.

This paper argues that queue-worker task profiling should be treated as a
task-kind-specific behavior modeling problem rather than as a problem of
collecting more raw telemetry. The central claim is that task attempts should
be profiled through a model that separates task identity, intrinsic demand,
execution context, resource demand, worker occupancy, uncertainty, and
operational risk. Within this framing, the paper makes the following
contributions.

First, we define a task cost model for queue workers that separates intrinsic
demand, contextual amplification, additive overhead, worker occupancy, and
residual uncertainty. This model treats task cost as conditional behavior
rather than as a fixed scalar attached to a task kind, queue, or
payload-size bucket.

Second, we introduce influence domains as a semantic decomposition of
queue-worker task behavior. These domains distinguish task identity, intrinsic
demand, execution context, resource demand, worker occupancy, outcomes,
uncertainty, risk, and profiling observations. The purpose of this
decomposition is to avoid flat metric lists, separate context factors from
symptoms and outcomes, and make explicit which kinds of observations are
relevant for a task kind.

Third, we propose Task Behavior Graphs as a task-kind- or
task-family-level representation that connects canonical features across
influence domains. A Task Behavior Graph describes how intrinsic task features
and execution-context factors contribute to resource demand, worker occupancy,
uncertainty, and operational risk. The graph provides the structure, while
task instances and attempts provide concrete feature values, context
observations, outcomes, and residuals for evaluation.

Fourth, we describe a metrics-to-features model based on feature registries,
metric descriptors, task behavior manifests, and domain packs. This mechanism
allows domain-specific raw observations to be mapped into canonical features
without requiring a globally complete metric list in the core profiling model.
It also makes missing or unexplained influences visible through residual error
and reduced confidence.

Fifth, we outline a profiling-policy selection flow that uses predictability,
context sensitivity, worker occupancy, operational consequences, and profiling
overhead to choose an appropriate observation strategy. Depending on the task
kind, attempt context, and risk profile, the selected strategy may range from
minimal accounting to detailed accounting, sampled tracing, tail-aware
profiling, diagnostic profiling, or isolated diagnostic execution.

The remainder of the paper is organized as follows.
Section~\ref{sec:terminology-notation} defines the terminology and notation
used throughout the paper. Sections~\ref{sec:background} and
\ref{sec:problem-statement} introduce the background and formulate the
profiling problem. Sections~\ref{sec:task-cost-model} through
\ref{sec:task-behavior-graphs} define the cost model, influence domains, task
taxonomy, and Task Behavior Graphs. Sections~\ref{sec:metrics-signals-features}
through~\ref{sec:profiling-policy-selection} describe observation mapping,
profiling methods, worker observation points, uncertainty, risk, and
profiling-policy selection. Section~\ref{sec:task-behavior-manifests}
introduces task behavior manifests and domain packs, and
Section~\ref{sec:case-studies} applies the model to representative task kinds.
The remaining sections discuss evaluation methodology, limitations, validity
threats, related work, and conclusions.